\documentclass{exam-zh}
\usepackage{mhchem}     % 化学式支持
\usepackage{siunitx}    % 物理单位支持

% 设置试题标题等信息
\title{模拟试卷}
\author{Exam-zh 示例}
\date{\today}

% 定义判断题勾叉符号（可选）
\newcommand{\correct}{\text{（正确）}}
\newcommand{\incorrect}{\text{（错误）}}

\begin{document}

% 生成试卷标题
\maketitle

% 第一部分：选择题
\section{选择题}

	\begin{question}
		$1+1$ 的值是多少？
		\begin{choices}
			\item 0
			\item 1
			\item 2
			\item 11
		\end{choices}
	\end{question}

	\begin{question}
		下面哪个是按顺序排列的英语大写字母表？
		\begin{choices}
			\item abcdefghijklmnopqrstuvwxyz
			\item ABCDEFGHIJKLMNOPQRSTUVWXYZ
			\item QWERTYUIOPASDFGHJKLZXCVBNM
			\item あいうえおかきくけこさしすせそたちつてと
		\end{choices}
	\end{question}

% 第二部分：填空题
\section{填空题}

\begin{question}
	配平下列化学方程式，并写出产物：\\
	\ce{CaCO3 + \_HCl -> } \fillin   % 留出6cm的填空线
	%\fillin[][6cm]   % 留出6cm的填空线
\end{question}

% 第三部分：判断题
\section{判断题}
	\begin{question}
		两直线平行，同旁内角相等。\hfill （ \quad 正确\quad 错误\quad ）
	\end{question}

% 第四部分：简答题
\section{简答题}
	\begin{question}
		一质点从原点静止出发，以加速度 $\SI{1}{\meter\per\square\second}$ 沿直线运动，求 $t = \SI{10}{\second}$ 时的位移大小。
		\vspace{4cm}   % 留出答题空间
	\end{question}

\end{document}